\section{Introduction}
% \textit{What is a hardware side channel, why are they bad, what is spectre in particular}
In a \textit{hardware side-channel attack} (HSCA), an unsafe instruction in a victim program creates operand-dependent hardware resource usage (e.g., load address operand-dependent cache hits/misses), and an attacker observes this resource usage (e.g., through its effect on program execution time) to infer the operand value. \textit{Constant-time} (CT) programming was long thought to offer a complete defense: it rewrites code to avoid passing secrets to unsafe instruction operands. However, \textit{Spectre attacks}\cite{spectre} (a class of HSCA) pass secrets to \textit{transient} (partially executed and eventually squashed due to hardware control-/data-flow mispredictions) unsafe instructions, revealing that secret program data can leak via hardware side channels even if it is never leaked or accessed architecturally. This means that even code that avoids passing any secrets to \textit{transmitters}, or instructions that leak one or more operands via an HSCA, can leak these secrets through transmitters under speculation. Spectre attacks circumvent pervasive software-level (e.g., CT programming or sandboxing) and hardware-level (e.g., privilege isolation) security policies, establishing HSCAs as a threat to all programs that hold secrets in architectural state.

% \textit{HSCDs exist for HSCAs }
Defending secret-processing programs against Spectre attacks is a prerequisite to running differently trusted code in the same process or on the same machine. Simply disabling the optimizations that enable these attacks incurs prohibitive performance overheads. Hence, numerous \textit{Spectre defenses} have been proposed that modify software, hardware, or both, with the goal of preventing speculative leakage of secrets. To prevent behavior from occurring speculatively without disabling optimization extensively, it is necessary to process information about speculation at runtime, so many Spectre defenses implement hardware mechanisms. We consider the class of defenses with some hardware component — such defenses can delay or taint-track speculative instructions until they are safe~\cite{yu2019stt,spt,nda,specshield,loughlin2021dolma,sakalis2019delay,recon,doppelganger,protean,cheng2024speclfb}, make speculative cache state invisible~\cite{yan2018invisispec,khasawneh2019safespec,ainsworth2020muontrap,ainsworth2021ghostminion}, undo state updates after a squash~\cite{saileshwar2019cleanupspec}, or execute speculative instructions data-obliviously~\cite{sdo}.

The performance implications of these hardware modifications are assessed via early-stage (pre-RTL) prototyping in computer architecture simulators, especially gem5\cite{gem5}. These prototypes aim to maximize performance while enforcing some advertised security guarantee \anna{should we use}. However, due to the largely inherent tradeoff between performance and security, subtle bugs that undermine a defense’s security may lead to severe underestimates of its performance overhead relative to an insecure baseline, rendering a prototype useless for predicting defense viability. 
Such bugs are common in state-of-the-art Spectre defenses: testing four simulated defenses from 2018 to 2024 (InvisiSpec, CleanupSpec, STT and SpecLFB), AMuLeT uncovered three known and six previously unknown bugs and vulnerabilities within three hours of testing~\cite{amulet}\cite{nickscommittoSTT}.

% not cycle accurate

Simulated Spectre defenses in hardware must be “correct enough” to yield reliable performance estimates to avoid wasted RTL development effort on nonviable proposals. At this stage of the design process, correctness is not formally proven but rather validated through testing, as formal proof would require a large investment of time and resources and designs are subject to change. However, \textit{security} correctness in Spectre defenses is a \textit{microarchitectural} property requiring knowledge of side channels, rather than a purely \textit{architectural} function, like functional correctness. While functional correctness can be reasonably validated by running benchmarks and checking their architectural outputs, it is more difficult to rigorously validate side-channel security guarantees \cite{spec2017}. There are three main reasons for this. First, there is no established formal framework for specifying the security properties of a defense within the space of possible Spectre attacks, nor terminology to specify how a violation occurs. Second, the existing concept of a litmus test is a short program, which by its nature can only specify architectural behavior, so is unable to consistently induce potential violations. Third, observing violating microarchitectural leakage requires probing microarchitectural state, which is not straightforward. 
% \todo{cite the spectre benchmarks}

% First, there is a lack of formal specification to describe the security guarantees a particular defense should uphold, making it difficult to test for this property. Second, there is no good suite of tests that exercises these security guarantees. \anna{is this a weird framing, I'm trying to set up the need for formal description} \anna{i need a consistent term for the overall security guarantee contract thing}

% \anna{it's important to note this is early in the design process => validation not verification, want to iterate, so we are working with test suites and checking, not with formal proofs}
There exists no principled way to validate security guarantees of defenses at this stage, which is necessary to do fast high-level iteration.
Existing validation either tests fabricated CPUs against generic speculation contracts~\cite{revizor,oleksenko2023hideseekspectresefficient}, verifies or fuzzes RTL designs~\cite{fadiheh2023exhaustive,contractshadow,rtl2upath,hur2022specdoctor,ghaniyoun2021introspectre}, or fuzzes simulator prototypes against a sequential baseline~\cite{amulet}; none checks a prototype against the defense's own security guarantee.
\notebr
\reviewnote{\textbf{Citation check} (asked for in an earlier source comment here): every key in this sentence and in the benchmarks sentence below exists in references.bib and is the intended paper: Revizor (ASPLOS'22), Hide and Seek with Spectres (S\&P'23), Fadiheh et al.'s exhaustive RTL detection (IEEE TC'23), Contract Shadow Logic (ASPLOS'25), RTL2M$\mu$PATH (MICRO'24), SpecDoctor (CCS'22), IntroSpectre (ISCA'21), AMuLeT (ASPLOS'25), Kocher's 15 Spectre v1 examples (2018) and Canella et al.\ (USENIX Security'19). AMuLeT's abstract confirms ``3 known and 6 unknown bugs and vulnerabilities, within 3 hours of testing'' (fu-2025-amulet.pdf p1). Three small points: (1) Pensieve~\cite{pensieve} (ISCA'23), which model-checks early-stage defense designs against speculative non-interference (yang-2023-pensieve.pdf abstract), is missing here and is the closest prior work for ``at this stage''; (2) \cite{spec2017} in the previous paragraph follows the security sentence but supports ``running benchmarks and checking their architectural outputs''; (3) in the second paragraph, \cite{nickscommittoSTT} (Mosier's STT commit 18e304f) follows the AMuLeT count, but that bug is not one of AMuLeT's nine; give it its own clause if it is meant as another example.}
\notebr
% When generating a set of litmus tests to validate some
% - we need litmus test suite
% - security different than functional correctness, because edge case could be very niche and that is very bad (can we make this argument), hard to randomly get to instantiation of one side channel but once the side channel is found it is infinitely exploitable
% - people have done this with randomized test generation, but it does not target potentially violating tests
% - people have benchmarks of tests that trigger spectre attacks, but they are not specialized to the 
Benchmarks of tests that trigger certain Spectre attacks have been compiled~\cite{kocher2018msvc,canella2019systematic}%
, but they are not specialized to the security guarantees of a particular defense, as different defenses enforce different rules on what hardware can leak, and cover a small subset of all relevant leakage scenarios. 
One piece of prior work — AMuLeT, has tested simulator implementations of Spectre defenses. AMuLeT uses randomized differentiated testing (running two randomly generated architecturally identical tests and checking for discrepancies in a limited portion of microarchitectural state) \cite{amulet}. 
\urgentbr
\urgentnote{The parenthetical misdescribes AMuLeT: it runs one random program on two inputs whose contract traces match and flags a difference in the microarchitectural trace (its Def.~2.1, fu-2025-amulet.pdf \S2.3), not ``two \ldots{} architecturally identical tests''. \textbf{PROPOSED TEXT} (edits the sentence just before this note): AMuLeT uses random relational testing: it runs a randomly generated program on pairs of inputs whose contract traces match and checks for differences in a limited part of the microarchitectural state~\cite{amulet}.}
\notebr
\anna{I think maybe drop the security property vs leakage predicate discussion here, and pick that up in related work. the main point being that it is random}.
% AMuLeT checks that a simulated defense satisfies a hardware--software contract~\cite{guarnieri2021contracts}: any two runs of a program whose contract traces match must also produce matching microarchitectural traces~\cite{amulet}. Such contracts generalize speculative non-interference (SNI)~\cite{guarnieri2020spectector}, which requires that speculatively executed instructions leak no more into microarchitectural state than non-speculative execution does. SNI and contracts are end-to-end, relational properties: they compare pairs of whole executions and do not say which instruction leaked or which rule of the defense it broke, so every violation must be root-caused by hand.
\anna{this paragraph abt fuzzing vs targeting? also in related work?}
% It also does not bias for generation of programs that are able to induce leakage. Security testing is sensitive, because just one leaky edge case that might be very difficult to randomly instantiate could be a potent attack vector, so it is important to generate a test suite that has a wide breadth of leaky instances. There exists no existing work that can generate targeted test sets for Spectre defense validation in a principled manner. 


In this paper we propose \toolname, the first rigorous, targeted, and automated approach for validating early-stage prototypes of hardware support for Spectre defenses against their promised security guarantees. \toolname provides (1) a \textbf{formal language} for specifying security guarantees of an arbitrary Spectre defense, (2) a \textbf{test generator} for generating a comprehensive set of tests (given a bound on test size) capturing all leakage scenarios blocked by a specified defense, and (3) a \textbf{test validator} harness for running these tests on gem5 implementations of the defense to detect bugs. \anna{stt running mini example figure 1?? bring the figure upwards otherwise}
\input{tikz/intro_threelayer_v2_8pt}% Fig. 1 right after the proposal (your note: bring the figure upwards), 2026-11-03

% \toolname enables validation of the hardware implementation Spectre defense against a defense security guarantee, and the specification of HSCAs against which it protects. This specification is precise to the cases blocked by a defense, and can be used to generate test sets; using it to prove adherence to a higher-level security property is future work.
% \notebr
\reviewnote{\textbf{ADD} ``(Fig.~\ref{fig:intro})'' to the sentence above: Fig.~1 is never referenced in the text, and this sentence is what it shows.}

\toolname{} tests defense implementations to their guarantees at the level of the defense's \emph{leakage predicate}: a condition on a single execution that names the transmitter, the protected operand it may not expose, and the speculative microarchitectural state in which the defense forbids the transmission. 

This means that an arbitrary natural language security description can be translated into formal components, which can then be used to generate all possible violating sequences of (abstract) instructions. These sequences are then used to generate a litmus test suite of assembly programs and corresponding \textit{microarchitectural annotations} that together induce the violating sequences. Finally, \toolname provides a harness that instruments the simulator to follow these speculation annotations, inducing the modeled violations.

The test sets generated by \toolname are complete in including all  \textit{leakage backbones}, or fundamental leakage structures, up to an instruction bound, and are diversified across timing, instruction type, and speculation direction during their concretization to assembly programs. 
\urgentbr
\urgentnote{Completeness holds only in part. Stage~1a (minimal graphs, up to six instructions) runs to UNSAT for STT and ReCon and their store-bypass variants, but SPT's Stage~1a does not converge (its four-instruction partition hits the 200,000 cap), and Stage~1b and Stage~2 are capped or sampled by design. \textbf{PROPOSED TEXT} (edits the sentence just before this note): For STT and ReCon, the test sets include every \textit{leakage backbone} up to six instructions (for SPT, up to a cap), and are diversified \ldots{} \textit{Source:} simspect wiki/state/EVAL\_MANIFEST.md `Stage 1' and `Stage 1.5' sections.}
\notebr
The harness that they run on both enforces the microarchitectural constraints of the test execution, and induces microarchitectural behavior that is likely to precipitate the test. 

In summary, \toolname provides not only a simple formal language in which to define a defense, but also uses it to target \textit{all possible} types of violations stopped by the defense, probing leakage in a principled and efficient manner and providing better assurance of design correctness.

We have expressed 10 hardware defenses and their variants in the \toolname{} language; STT, SpecShield, NDA, DoM, SDO, SPT, ReCon, InvisiSpec, SpecLFB, and CleanupSpec. We can therefore generate comprehensive litmus test suites up to specified bounds for each defense. 
\urgentbr
\urgentnote{Not supported as written. Only STT, ReCon and SPT exist as Alloy models (STT and ReCon also in store-bypass variants); InvisiSpec, CleanupSpec and SpecLFB are tested with STT's test set and a leakage criterion written from each defense's paper; SpecShield, NDA, DoM and SDO have no model. \textbf{PROPOSED TEXT} (edits the sentence just before this note): We have expressed three hardware defenses in the \toolname{} language, STT, ReCon and SPT, and test three more, InvisiSpec, CleanupSpec and SpecLFB, with STT's test set and a defense-specific leakage criterion. \textit{Source:} simspect STAGE1\_alloy/models/ (STT\_6, Recon\_6, SPT\_6\_oneLP\_v2 and their proposed/ultimate versions); specs/leakage/\{invisispec,cleanupspec,speclfb\}.json.}
\notebr

% We have run the \toolname{} testing pipeline with an instruction bound of six on the gem5 implementations of STT, SPT, and ReCon, and on the AMuLeT builds of InvisiSpec, SpecLFB, and CleanupSpec

We apply \toolname{}, with an instruction bound of six, to every simulator prototype of a state-of-the-art defense we could obtain: the gem5 implementations of STT, SPT, and ReCon, and the AMuLeT builds of InvisiSpec, CleanupSpec, and SpecLFB. Checking each defense against only the leakage its own paper claims to prevent, \toolname{} replicates every previously known bug in that scope: InvisiSpec's speculative L1D eviction, CleanupSpec's three cleanup bugs (an uncleaned speculative read-modify-write, an uncleaned line-crossing access, and cleaning too much), and SpecLFB's unprotected first load.
\notebr
\reviewnote{Evidence behind this sentence, for the counts you add later: AMuLeT reports nine bugs on these defenses; five fall inside each defense's own guarantee and SimSpect finds all five: UV1 (InvisiSpec, 2,995 of 3,000 tests), UV3 (only with the read-modify-write corpus and a cold cache: 407 of 590 tests cold, 0 of 590 warm), UV4 (38 tests), UV5 (one test so far, inst-039102; whether to count it is still your open decision) and UV6 (about 1,140 of 1,500 tests per mode). Out of scope, each by its paper: UV2 (MSHR contention, InvisiSpec \S VII), KV1 (instruction fetch, InvisiSpec fn.~2), KV2 (unXpec cleanup timing: CleanupSpec \S4(b) p8 acknowledges that cleanup time varies and leaves constant-time cleanup to future work, and AMuLeT observes it through the L1I, which CleanupSpec \S2.1 p3 excludes) and KV3 (speculative store fills the D-TLB: STT assumes stores act only at retirement, \S5 p5). \textit{Source:} simspect wiki/findings/mdsToRead/AMULET\_BUGS\_DETECTION\_ACCOUNTING.md sec.~2; URGENT\_TODO.md (KV2, KV3, UV3, UV5 items).}
\notebr

\toolname{} also finds six previously unknown bugs. (1) STT fails to taint one register: a check that skips register index 16 ignores the register class, so a secret copied through the low half of \texttt{xmm4} loses its taint. In ReCon's underlying implementation of STT, (2) a load that receives its value by store-to-load forwarding completes before the taint check runs, and (3) an instruction that merges two tainted operands keeps only the oldest taint source, so its result becomes untainted once that source's branch resolves. SpecLFB (4) misimplements its ROB walk order, so loads after the first can also go unprotected, and (5) fails to protect loads in the shadow of a direct \texttt{jmp} or \texttt{call}, which erases an older unresolved branch's shadow. Finally, (6) STT and SPT share a visibility-point/squash timing race: a wrong-path transmitter is released when its branch resolves and accesses the cache before the squash removes it.
\urgentbr
\urgentnote{(6) leaves out the condition the race needs: it leaks only when commit's squash reaches the issue stage at least two to three cycles late (gem5's commitToIEWDelay), and never at stock timing. SPT Fence on 24 race tests: stock 0, delay 2 $\to$ 3, delay 3 $\to$ 14; STT leaks at delay 3 and 0 of 1,200 at stock. This is a global latency setting, not one of the test hooks, so say so. \textbf{PROPOSED TEXT} (edits the last sentence before this note): Finally, (6) STT and SPT share a visibility-point/squash timing race: when the squash signal reaches the issue stage a few cycles after commit (a longer commit-to-IEW delay than gem5's default), a wrong-path transmitter is released when its branch resolves and accesses the cache before the squash removes it. Also, (5) is reached by the generated tests' own jumps, not only by the harness's call. \textit{Source:} simspect wiki/findings/bug/BUG\_TAXONOMY.md A4, A5; project log 2026-11-09 15:20 (bug-investigate-a5) and 20:10 (bug-investigate-e14new).}
\notebr

Compared to prior work, \toolname{} needs [N]$\times$ fewer tests on average to find the known bugs. Every test it generates contains a protected transmitter, while only 17\% of AMuLeT's randomly synthesized tests contain the components necessary for leakage.

% NO

% Because the leakage predicate is formal, it can be related to the end-to-end properties. A bounded Alloy check that every pair of executions allowed by a defense's leakage predicate satisfies SNI, or a given contract, would turn passing \toolname{} tests into evidence for that end-to-end claim, such as STT's non-interference theorem~\cite{yu2019stt}.